Badamy problem rozszerzenia zadanej bijekcji pięciu punktów płaszczyzny $\Q^2$ do automorfizmu
wielomianowego o współczynnikach wymiernych i możliwie najmniejszym stopniu; ogólna konstrukcja
daje stopień $(m-1)^2=16$. Najniższy stopień rozwiązania przyjętego przez niezależny weryfikator
wynosi $6$ (złożenie dwóch odwzorowań trójkątnych stopni 2 i 3 z czynnikami afinicznymi).
Korzystając z cytowanego twierdzenia Junga--van der Kulka pokazujemy, że stopień każdego
rozwiązania jest iloczynem co najmniej dwóch liczb $\ge2$, a dwa fakty obliczeniowe (brak
niestałego funkcjonału afinicznego zgodnego na danych oraz brak wymiernych rozwiązań o
multistopniu $(2,2)$, potwierdzony dwiema metodami) wykluczają stopnie poniżej 6.
Zatem, modulo cytowane twierdzenie i rachunki komputerowe, najmniejszy stopień wynosi dokładnie 6;
nad $\Q(\sqrt{94})$ i pewnym ciałem kubicznym stopień 4 byłby osiągalny.
Otwarte pozostają m.in. ręczny dowód wykluczenia stopnia 4, opis wszystkich rozwiązań stopnia 6
oraz zachowanie problemu dla danych ogólnych i większych $m$.
